% Seminar 2 - Page 6. Fragment for % Seminar 2 - Page 6. Fragment for % Seminar 2 - Page 6. Fragment for % Seminar 2 - Page 6. Fragment for \input{seminars/pages/seminar2-pagina6-en}.
% No preamble and no forced page break.
% EDITORIAL NOTES:
% - This page continues the argument for Problem 7 from Page 5.
% - Several intermediate expressions are crossed out or overwritten;
%   only legible, uncancelled expressions appear in the clean text.
% - The indices u_1,u_2,u_3 in the descending sequence follow the
%   handwriting; the repeated construction is abbreviated in the source.
% - In Problem 8, Hom(A,B) denotes all maps A -> B, not explicitly
%   homomorphisms of an algebraic structure.
% - Problem 9 asserts there is no surjection X -> P(X) without a proof
%   on this manuscript page.
% - Problem 10 proves only injectivity of phi; a proof of surjectivity
%   has not been invented or added.

\subsection*{Seminar 2 - Page 6}

\noindent\underline{Solution 7 (continued).}
\[
\Longrightarrow f(n)\geq k+1.
\]
\noindent Hence
\[
 f(k)\geq k\qquad (\forall k\in\mathbb{N}^{*}).
\]

\noindent Suppose, for a contradiction, that there exists
$k\in\mathbb{N}^{*}$ such that
\[
 f(k)>k.
\]
\noindent We show that if $n>k$, then $f(n-1)>k$.
\[
 n>k\ \Longrightarrow\ n\geq k+1\ \Longrightarrow\ n-1\geq k
 \ \Longrightarrow\ f(n-1)\geq k.
\]
\noindent Suppose $f(n-1)=k$. Then
\[
 k=f(n-1)\geq n-1\geq k
 \quad\Longrightarrow\quad n-1=k
 \quad\Longrightarrow\quad f(k)=k,\quad\text{a contradiction}.
\]

\noindent Let $u_1>k$. Then
\[
 f(u_1)>f(u_2),\qquad u_2=f(u_1-1)>k.
\]
\noindent Similarly,
\[
 f(u_2)>f(u_3),\qquad u_3=f(u_2-1)>k.
\]
\noindent Continuing indefinitely, we obtain an infinite strictly
 decreasing sequence of natural numbers.

\bigskip
\noindent\fbox{8}\quad
\[
 \operatorname{Hom}(A,B)\overset{\mathrm{def}}{=}
 \{f\mid f:A\longrightarrow B\}.
\]
\noindent $\mathcal{P}(X)$ is in bijection with
\[
 \operatorname{Hom}(X,\{0,1\})
   =\{\varphi\mid\varphi:X\longrightarrow\{0,1\}\}.
\]

\medskip
\noindent\fbox{9}\quad There is no surjective map
\[
 f:X\longrightarrow\mathcal{P}(X).
\]

\medskip
\noindent\fbox{10}\quad
\[
 \varphi:\mathcal{P}(X)\longrightarrow
 \operatorname{Hom}(X,\{0,1\}).
\]
\noindent If $A\subseteq X$, let
$\varphi(A)=\varphi_A:X\longrightarrow\{0,1\}$, where
\[
 \varphi_A(x)=
 \begin{cases}
  1,&x\in A,\\
  0,&x\notin A.
 \end{cases}
\]
\noindent $\varphi$ is injective: if
\[
 \varphi(A)=\varphi(B),\qquad \varphi_A=\varphi_B,
\]
\noindent then
\[
 x\in A\ \Longleftrightarrow\ \varphi_A(x)=1
 \ \Longleftrightarrow\ \varphi_B(x)=1
 \ \Longleftrightarrow\ x\in B
 \quad\Longrightarrow\quad A=B.
\]
.
% No preamble and no forced page break.
% EDITORIAL NOTES:
% - This page continues the argument for Problem 7 from Page 5.
% - Several intermediate expressions are crossed out or overwritten;
%   only legible, uncancelled expressions appear in the clean text.
% - The indices u_1,u_2,u_3 in the descending sequence follow the
%   handwriting; the repeated construction is abbreviated in the source.
% - In Problem 8, Hom(A,B) denotes all maps A -> B, not explicitly
%   homomorphisms of an algebraic structure.
% - Problem 9 asserts there is no surjection X -> P(X) without a proof
%   on this manuscript page.
% - Problem 10 proves only injectivity of phi; a proof of surjectivity
%   has not been invented or added.

\subsection*{Seminar 2 - Page 6}

\noindent\underline{Solution 7 (continued).}
\[
\Longrightarrow f(n)\geq k+1.
\]
\noindent Hence
\[
 f(k)\geq k\qquad (\forall k\in\mathbb{N}^{*}).
\]

\noindent Suppose, for a contradiction, that there exists
$k\in\mathbb{N}^{*}$ such that
\[
 f(k)>k.
\]
\noindent We show that if $n>k$, then $f(n-1)>k$.
\[
 n>k\ \Longrightarrow\ n\geq k+1\ \Longrightarrow\ n-1\geq k
 \ \Longrightarrow\ f(n-1)\geq k.
\]
\noindent Suppose $f(n-1)=k$. Then
\[
 k=f(n-1)\geq n-1\geq k
 \quad\Longrightarrow\quad n-1=k
 \quad\Longrightarrow\quad f(k)=k,\quad\text{a contradiction}.
\]

\noindent Let $u_1>k$. Then
\[
 f(u_1)>f(u_2),\qquad u_2=f(u_1-1)>k.
\]
\noindent Similarly,
\[
 f(u_2)>f(u_3),\qquad u_3=f(u_2-1)>k.
\]
\noindent Continuing indefinitely, we obtain an infinite strictly
 decreasing sequence of natural numbers.

\bigskip
\noindent\fbox{8}\quad
\[
 \operatorname{Hom}(A,B)\overset{\mathrm{def}}{=}
 \{f\mid f:A\longrightarrow B\}.
\]
\noindent $\mathcal{P}(X)$ is in bijection with
\[
 \operatorname{Hom}(X,\{0,1\})
   =\{\varphi\mid\varphi:X\longrightarrow\{0,1\}\}.
\]

\medskip
\noindent\fbox{9}\quad There is no surjective map
\[
 f:X\longrightarrow\mathcal{P}(X).
\]

\medskip
\noindent\fbox{10}\quad
\[
 \varphi:\mathcal{P}(X)\longrightarrow
 \operatorname{Hom}(X,\{0,1\}).
\]
\noindent If $A\subseteq X$, let
$\varphi(A)=\varphi_A:X\longrightarrow\{0,1\}$, where
\[
 \varphi_A(x)=
 \begin{cases}
  1,&x\in A,\\
  0,&x\notin A.
 \end{cases}
\]
\noindent $\varphi$ is injective: if
\[
 \varphi(A)=\varphi(B),\qquad \varphi_A=\varphi_B,
\]
\noindent then
\[
 x\in A\ \Longleftrightarrow\ \varphi_A(x)=1
 \ \Longleftrightarrow\ \varphi_B(x)=1
 \ \Longleftrightarrow\ x\in B
 \quad\Longrightarrow\quad A=B.
\]
.
% No preamble and no forced page break.
% EDITORIAL NOTES:
% - This page continues the argument for Problem 7 from Page 5.
% - Several intermediate expressions are crossed out or overwritten;
%   only legible, uncancelled expressions appear in the clean text.
% - The indices u_1,u_2,u_3 in the descending sequence follow the
%   handwriting; the repeated construction is abbreviated in the source.
% - In Problem 8, Hom(A,B) denotes all maps A -> B, not explicitly
%   homomorphisms of an algebraic structure.
% - Problem 9 asserts there is no surjection X -> P(X) without a proof
%   on this manuscript page.
% - Problem 10 proves only injectivity of phi; a proof of surjectivity
%   has not been invented or added.

\subsection*{Seminar 2 - Page 6}

\noindent\underline{Solution 7 (continued).}
\[
\Longrightarrow f(n)\geq k+1.
\]
\noindent Hence
\[
 f(k)\geq k\qquad (\forall k\in\mathbb{N}^{*}).
\]

\noindent Suppose, for a contradiction, that there exists
$k\in\mathbb{N}^{*}$ such that
\[
 f(k)>k.
\]
\noindent We show that if $n>k$, then $f(n-1)>k$.
\[
 n>k\ \Longrightarrow\ n\geq k+1\ \Longrightarrow\ n-1\geq k
 \ \Longrightarrow\ f(n-1)\geq k.
\]
\noindent Suppose $f(n-1)=k$. Then
\[
 k=f(n-1)\geq n-1\geq k
 \quad\Longrightarrow\quad n-1=k
 \quad\Longrightarrow\quad f(k)=k,\quad\text{a contradiction}.
\]

\noindent Let $u_1>k$. Then
\[
 f(u_1)>f(u_2),\qquad u_2=f(u_1-1)>k.
\]
\noindent Similarly,
\[
 f(u_2)>f(u_3),\qquad u_3=f(u_2-1)>k.
\]
\noindent Continuing indefinitely, we obtain an infinite strictly
 decreasing sequence of natural numbers.

\bigskip
\noindent\fbox{8}\quad
\[
 \operatorname{Hom}(A,B)\overset{\mathrm{def}}{=}
 \{f\mid f:A\longrightarrow B\}.
\]
\noindent $\mathcal{P}(X)$ is in bijection with
\[
 \operatorname{Hom}(X,\{0,1\})
   =\{\varphi\mid\varphi:X\longrightarrow\{0,1\}\}.
\]

\medskip
\noindent\fbox{9}\quad There is no surjective map
\[
 f:X\longrightarrow\mathcal{P}(X).
\]

\medskip
\noindent\fbox{10}\quad
\[
 \varphi:\mathcal{P}(X)\longrightarrow
 \operatorname{Hom}(X,\{0,1\}).
\]
\noindent If $A\subseteq X$, let
$\varphi(A)=\varphi_A:X\longrightarrow\{0,1\}$, where
\[
 \varphi_A(x)=
 \begin{cases}
  1,&x\in A,\\
  0,&x\notin A.
 \end{cases}
\]
\noindent $\varphi$ is injective: if
\[
 \varphi(A)=\varphi(B),\qquad \varphi_A=\varphi_B,
\]
\noindent then
\[
 x\in A\ \Longleftrightarrow\ \varphi_A(x)=1
 \ \Longleftrightarrow\ \varphi_B(x)=1
 \ \Longleftrightarrow\ x\in B
 \quad\Longrightarrow\quad A=B.
\]
.
% No preamble and no forced page break.
% EDITORIAL NOTES:
% - This page continues the argument for Problem 7 from Page 5.
% - Several intermediate expressions are crossed out or overwritten;
%   only legible, uncancelled expressions appear in the clean text.
% - The indices u_1,u_2,u_3 in the descending sequence follow the
%   handwriting; the repeated construction is abbreviated in the source.
% - In Problem 8, Hom(A,B) denotes all maps A -> B, not explicitly
%   homomorphisms of an algebraic structure.
% - Problem 9 asserts there is no surjection X -> P(X) without a proof
%   on this manuscript page.
% - Problem 10 proves only injectivity of phi; a proof of surjectivity
%   has not been invented or added.

\subsection*{Seminar 2 - Page 6}

\noindent\underline{Solution 7 (continued).}
\[
\Longrightarrow f(n)\geq k+1.
\]
\noindent Hence
\[
 f(k)\geq k\qquad (\forall k\in\mathbb{N}^{*}).
\]

\noindent Suppose, for a contradiction, that there exists
$k\in\mathbb{N}^{*}$ such that
\[
 f(k)>k.
\]
\noindent We show that if $n>k$, then $f(n-1)>k$.
\[
 n>k\ \Longrightarrow\ n\geq k+1\ \Longrightarrow\ n-1\geq k
 \ \Longrightarrow\ f(n-1)\geq k.
\]
\noindent Suppose $f(n-1)=k$. Then
\[
 k=f(n-1)\geq n-1\geq k
 \quad\Longrightarrow\quad n-1=k
 \quad\Longrightarrow\quad f(k)=k,\quad\text{a contradiction}.
\]

\noindent Let $u_1>k$. Then
\[
 f(u_1)>f(u_2),\qquad u_2=f(u_1-1)>k.
\]
\noindent Similarly,
\[
 f(u_2)>f(u_3),\qquad u_3=f(u_2-1)>k.
\]
\noindent Continuing indefinitely, we obtain an infinite strictly
 decreasing sequence of natural numbers.

\bigskip
\noindent\fbox{8}\quad
\[
 \operatorname{Hom}(A,B)\overset{\mathrm{def}}{=}
 \{f\mid f:A\longrightarrow B\}.
\]
\noindent $\mathcal{P}(X)$ is in bijection with
\[
 \operatorname{Hom}(X,\{0,1\})
   =\{\varphi\mid\varphi:X\longrightarrow\{0,1\}\}.
\]

\medskip
\noindent\fbox{9}\quad There is no surjective map
\[
 f:X\longrightarrow\mathcal{P}(X).
\]

\medskip
\noindent\fbox{10}\quad
\[
 \varphi:\mathcal{P}(X)\longrightarrow
 \operatorname{Hom}(X,\{0,1\}).
\]
\noindent If $A\subseteq X$, let
$\varphi(A)=\varphi_A:X\longrightarrow\{0,1\}$, where
\[
 \varphi_A(x)=
 \begin{cases}
  1,&x\in A,\\
  0,&x\notin A.
 \end{cases}
\]
\noindent $\varphi$ is injective: if
\[
 \varphi(A)=\varphi(B),\qquad \varphi_A=\varphi_B,
\]
\noindent then
\[
 x\in A\ \Longleftrightarrow\ \varphi_A(x)=1
 \ \Longleftrightarrow\ \varphi_B(x)=1
 \ \Longleftrightarrow\ x\in B
 \quad\Longrightarrow\quad A=B.
\]
